\section{The Bill for a Shareable \code{node}}
\label{sec:ch09-the-bill-for-a-shareable-node}

\index{node@\code{node(id:)}!across a seam|(}%
Chapter~\ref{ch:composition} spent thirty-three lines on a type interceptor so
that \code{node} and \code{nodes} would carry \code{@shareable} and two
subgraphs would compose. I said at the time that the trade was worth it,
because the alternative promised four fields in order to buy a promise about
two. There is now a second subgraph, and the promise about two comes due.

\index{shareable@\code{"@shareable}!says any subgraph can answer}%
\code{@shareable} means any subgraph declaring the field can resolve it. Read
that against what the two services can actually do. Sessions resolves
\code{node(id:)} for a \code{Session}. Speakers resolves it for a
\code{Speaker}. Neither can do the other's, and neither says so anywhere,
because there is nowhere in the schema to say it.

\index{node@\code{node(id:)}!works when the query names the type}%
Ask the router for a speaker by node id, naming the type in a fragment, and
everything is fine:

\begin{minted}{graphql}
{ node(id: "U3BlYWtlcjox") { __typename ... on Speaker { name } } }
\end{minted}

\begin{minted}{json}
{"data":{"node":{"__typename":"Speaker","name":"Ada Fischer"}}}
\end{minted}

\index{node@\code{node(id:)}!the fragment is what routes it}%
The fragment is what made that work. It is the only thing in the request that
names a type, and the query went to the service that owns that type; take it
away, as the next request does, and the query goes somewhere else.

\index{node@\code{node(id:)}!the same id with the wrong fragment fails}%
Now the same identifier and the same field, with a fragment for the other type:

\begin{minted}{graphql}
{ node(id: "U3BlYWtlcjox") { __typename ... on Session { title } } }
\end{minted}

\begin{minted}[fontsize=\footnotesize,breakanywhere]{json}
{"errors":[{"message":"Failed to fetch from Subgraph 'sessions'.","extensions":{"errors":[{"message":"There is no node resolver registered for type `Speaker`.","path":["node"],"extensions":{"code":"DOWNSTREAM_SERVICE_ERROR"}}],"serviceName":"sessions","statusCode":200}}],"data":{"node":null}}
\end{minted}

\index{node@\code{node(id:)}!chapter 4's rule does not survive the seam}%
Chapter~\ref{ch:schema-design} established what this request does in a single
service, and made a point of it: a well-formed node id of the wrong type is
\emph{not} an error. The fragment does not match, the client gets an object
with no fields it asked for, and nothing anywhere is wrong. That was a nice
property and it does not survive a seam. With no \code{Speaker} named anywhere
in the request, the planner sent it to Sessions, which does not know what a
\code{Speaker} is. Why Sessions rather than Speakers I cannot tell you; both
declare \code{node} and the schema gives the planner nothing to prefer one
by.

\index{nodes@\code{nodes(ids:)}!half a list}%
\code{nodes} makes the same point with the failure in the middle rather than at
the end. Two identifiers, one of each type, which no single subgraph can
answer:

\begin{minted}{graphql}
{ nodes(ids: ["U2Vzc2lvbjox", "U3BlYWtlcjox"]) { __typename } }
\end{minted}

\begin{minted}[fontsize=\footnotesize,breakanywhere]{json}
{"data":{"nodes":[{"__typename":"Session"},null]},"errors":[{"message":"Failed to fetch from Subgraph 'sessions'."}]}
\end{minted}

\index{nodes@\code{nodes(ids:)}!answers 200 with half the answer}%
HTTP 200, one type resolved, one null, one error. The list is the shape the
client asked for and half of it is missing.

\subsection{Whose Fault This Is}

\index{node@\code{node(id:)}!is a routing problem, not a resolver one}%
Nothing in either service is broken. Ask each of them for its own
identifiers on its own port and both answer correctly, and the composed
\code{Node} interface at the router names both types, so the schema is not
lying about what exists. The router cannot tell, from an opaque
identifier and no fragment, which of two services to ask.

\index{node id!opacity has a cost after all}%
Which is the bill for something chapter~\ref{ch:schema-design} was pleased
about. That chapter said the node id format is not a secret, decoded one on the
page, and warned against reading opacity as a control. This is the other half:
the format is not a secret from you, and it \emph{is} a secret from the router,
because a router that decoded one would be depending on an encoding no
specification requires.

\index{node@\code{node(id:)}!four fixes, and none of them is good}%
I looked for a fix and did not find one worth shipping. Four things you might
try:

\begin{itemize}
  \item Give every subgraph a node resolver for every type. They do not have
        the rows. Doing it anyway would be undoing the split.
  \item Declare \code{node} in one subgraph only, so nothing collides and no
        \code{@shareable} is needed. Composition passes and the failure is
        identical, because that one subgraph still owns one type out of two.
  \item Have the node resolver return null for a type it does not own rather
        than throwing. That trades a loud failure for a quiet wrong answer,
        which is the class of defect chapters~\ref{ch:representation-order}
        and~\ref{ch:entities-authorization} exist to warn about.
  \item Put \code{@inaccessible} on \code{node} and \code{nodes} everywhere and
        take the field out of the client schema. Honest, and it costs the
        reader the whole of chapter~\ref{ch:schema-design}'s refetch story.
\end{itemize}

\index{node id!what federating on one costs}%
So the graph ships as measured and I am telling you where the edge is rather
than papering over it. This book federates on the global node id because that
is what chapter~\ref{ch:schema-design} already had and no better key was
available, and I said in chapter~\ref{ch:federation-model} that the choice was
my judgment and not Apollo's guidance, because Apollo publishes none. This is
what the judgment costs. A client that always names the type it expects, which
is what a typed client generated from the schema does, never meets it. A client
that passes an identifier it was handed and asks what it is, meets it
immediately.

\index{satisfiability!the router can route only on what the query says}%
Chapter~\ref{ch:satisfiability} is where this gets picked up, because it is the
chapter about the composer deciding whether every field can be reached from a
root field, and \code{node} is the root field that makes that question hardest.
\index{node@\code{node(id:)}!across a seam|)}%
